% TOP_PROBLEM: {"id": 3167, "title": "r-regular graphs are not uniquely hamiltonian.", "category_id": 3, "category": {"id": 3, "name": "graph_theory", "display_name": "Graph Theory", "description": "Problems involving graphs, networks, and their properties.", "slug": "graph-theory", "order_index": 3, "created_at": "2026-07-31T15:26:25.671Z"}, "status": "open", "problem_number": "OPG-480", "set_id": 12, "statusQualification": "The graph is explicitly simple, as required by the source comments excluding multigraph counterexamples."}
% FIRST_POSTED: 2026-10-01
% ENTRY_KIND: statement_only
% UnsolvedMath ID 3167; source catalogue code OPG-480
% !TeX program = lualatex
% Definition and sources only; this file is not an LLM solution attempt.
\documentclass[11pt,a4paper]{article}
\usepackage[margin=25mm]{geometry}
\usepackage{fontspec}
\setmainfont{lmroman10-regular.otf}[BoldFont=lmroman10-bold.otf,
 ItalicFont=lmroman10-italic.otf,BoldItalicFont=lmroman10-bolditalic.otf]
\usepackage{amsmath,amssymb,mathtools}
\usepackage{unicode-math}
\setmathfont{latinmodern-math.otf}
\AtBeginDocument{\let\setminus\smallsetminus}
\usepackage{xurl,hyperref}
\hypersetup{colorlinks=true,urlcolor=blue}
\setcounter{secnumdepth}{0}
\setlength{\parindent}{0pt}
\setlength{\parskip}{5pt}
\setlength{\emergencystretch}{3em}
\begin{document}
\section*{r-regular graphs are not uniquely hamiltonian.}\label{3167}
{\small UnsolvedMath ID 3167 · OPG-480 · Graph Theory · OpenGarden}
\subsection{Definitions and mathematical statement}
Let $r\geq3$ be an integer and let $G=(V,E)$ be a finite simple
undirected $r$-regular graph: every vertex has exactly $r$
neighbours. A Hamilton cycle is a connected spanning
$2$-regular subgraph, equivalently a cycle passing through
all vertices once. Count Hamilton cycles by their edge sets.
Changing the starting vertex or reversing the direction of
traversal does not create a new cycle.

The graph is uniquely Hamiltonian when it has exactly one
Hamilton cycle. The conjecture asserts that no finite simple
$r$-regular graph is uniquely Hamiltonian for any $r>2$.
Writing $h(G)$ for the number of Hamilton cycles, its target
is
\[
 h(G)\neq1\qquad
 \text{for every finite simple regular graph of degree }r\geq3.
\]
Equivalently, whenever such a $G$ has a Hamilton cycle $C$,
there must be another Hamilton cycle $C'$ with
$E(C')\neq E(C)$.

The assertion permits graphs with no Hamilton cycle; it
does not assert that regularity alone guarantees one. The
two cycles may share vertices and edges and need not be
edge-disjoint. No connectivity assumption beyond that
forced by existence of the first Hamilton cycle is needed.
Simplicity is essential to the intended source question:
parallel edges are excluded, as emphasized in its discussion.
Degree two is omitted because a connected simple cycle is
itself a uniquely Hamiltonian $2$-regular graph.
\subsection{Short English statement}
Must every Hamiltonian regular simple graph of degree greater than two have at least two distinct Hamilton cycles?
\subsection{Sources}
\begin{itemize}
\item[S1] ULAM.AI and the UnsolvedMath Contributors, supplied problems.json, record 3167 (OPG-480), OpenGarden collection. Catalogue identity and source transcription; CC BY 4.0.
\url{https://huggingface.co/datasets/ulamai/UnsolvedMath}.
\item[S2] Open Problem Garden, r-regular graphs are not uniquely hamiltonian.. Original problem and discussion; the mathematical formulation below is expanded from the supplied source record.
\url{http://www.openproblemgarden.org/op/uniquely_hamiltonian_graphs}.
\end{itemize}
\end{document}
